\subsection{The precise reciprocity condition for canonical central lifts}

\begin{claim}{D.27}{Proved; an exact compatibility criterion}\label{clm:D27}
Let $k$ be a number field and, for each place $v$, prescribe a
compatible local class $\lambda_v\in\widehat{k_v^\times}$.
At real places this is a sign class, and at complex places it
is trivial. Write $\lambda_{v,n}\in H^1(k_v,\mu_n)$ for its
component. There exists a global
$A\in\widehat{k^\times}$ with these localizations if and only
if, for every $n\ge1$, the following conditions hold:
\begin{enumerate}
\item $\lambda_{v,n}$ is unramified at almost every finite
place $v$;
\item for every $\chi\in H^1(k,\mathbf Z/n)$,
\[
 \sum_v\operatorname{inv}_v
       (\lambda_{v,n}\smile\chi_v)=0,                       \tag{D.27.1}
\]
with the usual Tate convention at real places.
\end{enumerate}
Such an $A$ is unique. Reconstructions of prescribed local
data are therefore compatible with restriction, multiplication
and norms whenever those local data satisfy the corresponding
identities. The criterion applies componentwise to
quasi-trivial tori.

For $U=\mathbf P^1_k-\{0,1,\infty\}$,
$a\in k-\{0,1\}$ and $W=U-\{a\}$, the additional tests
for the canonical values $x_v-a$ are exactly the residue
summand in the split sequence
\[
 0\longrightarrow\operatorname{Br}(U)[n]
  \longrightarrow\operatorname{Br}(W)[n]
  \xrightarrow{\partial_a}H^1(k,\mathbf Z/n)
  \longrightarrow0.                                      \tag{D.27.2}
\]
Here $\partial_a$ is the cohomological residue, the negative
of the Witt residue under the Brauer identification. A
splitting is $\chi\mapsto\kappa(t-a)\smile\chi$.
Thus the existence of an arbitrary central lift in D.23 does
not impose the canonical local-value conditions, and D.26
does not assert their nonemptiness for a new interior puncture.
\end{claim}
\begin{proof}
Put $M=\mu_n$ and $M^D=\operatorname{Hom}(M,\mathbf G_m)
=\mathbf Z/n$. The first condition means exactly that
$\lambda_n=(\lambda_{v,n})_v$ belongs to the restricted
product $P^1(k,M)$ with respect to unramified cohomology.
Away from places dividing $n$, it says that the local
valuation is zero modulo $n$. For a fixed $\chi$, both
factors are unramified away from a finite set, and their
cup product is then zero: it comes from degree-two
cohomology of a finite residue field. Thus (D.27.1) is a
finite sum.

Poitou--Tate exactness gives
\[
 H^1(k,\mu_n)\longrightarrow P^1(k,\mu_n)
       \longrightarrow H^1(k,\mathbf Z/n)^*,              \tag{D.27.3}
\]
where the second map is the sum of the local cup products
\cite[I, \S4, Theorem 4.10(b)]{D_MilneADT2006}.
This is exactness for the \emph{actual global image}, not
only its closure. Indeed, the preceding Lemma 4.9 says that
localization is proper when global $H^1$ has the discrete
topology. The preimage of every compact-open local condition
is finite, by the bounded-degree, bounded-ramification
finiteness of number-field extensions. The global image
therefore meets each compact-open subset in a finite closed
set and is closed in $P^1$. The real local groups and the
restricted product used here are exactly the conventions
preceding Lemma 4.8 of that source.

The two conditions are consequently equivalent to existence
of a global preimage at each modulus. Let $S_n$ be its set
of global preimages. It is finite and nonempty: its difference
group is the finite localization kernel, by Lemma 4.9 or
Theorem 4.10(a). No injectivity of finite-level localization
is assumed. For $n\mid m$, reduction maps $S_m$ to $S_n$,
since the prescribed local components are compatible.
The inverse system of these finite nonempty sets has a
point. Explicitly, use the cofinal sequence $n=r!$;
any finite collection of compatibility conditions is realized
by an element of one sufficiently divisible level. These
closed conditions have the finite-intersection property in
the compact product of the finite sets $S_{r!}$.
This argument does not require the individual transition
maps to be surjective. The inverse-limit point is the
required completed class $A$.

If two completed classes have the given localizations,
their ratio is trivial at every finite place at every
modulus, so D.26 makes it trivial globally. This gives
uniqueness. Restriction, multiplication and the Kummer norm
map commute with localization; a local norm is the product
of the norms over the places above the given place.
The asserted functorialities now follow from uniqueness.
Shapiro's lemma applies this argument to
$\operatorname{Res}_{L/k}\mathbf G_m$, and products give
the quasi-trivial torus assertion.

To apply the criterion to central lifting, first choose a
global central lift, as D.23 permits, and specify the desired
local central lifts and their comparison paths. If $T_D$ is
the relevant quasi-trivial torus, their differences from the
chosen origin are classes in $H^1(k_v,T(T_D))$, cohomology
with coefficients in its full Tate module. The criterion above
is exactly the condition for one global difference class.
An arbitrary global lift supplies this origin; it does not
supply the desired local differences or prove the two
conditions for them.

We identify those conditions concretely for an added
rational puncture. In the inverted torus coordinate, the
canonical map from $W$ into the diagonal torus torsor has
coordinate $t-a$. At an ordinary $x_v\in W(k_v)$ its value
is $x_v-a$. This is a prescribed value, not an arbitrary
nonzero point in the torus fibre. If $x_v=a$, the value zero
is not a Kummer class: a new tangent class for the punctured
curve must be chosen. At an old cusp the evaluation is
instead determined using the already specified old
tangential section and the leading term of $t-a$. At a
real place the existing component does not silently select
a new sign or tangent. The criterion applies once these
local data have actually been specified.

For $\chi\in H^1(k,\mathbf Z/n)$ define
\[
 \beta_{a,\chi}=\kappa(t-a)\smile\chi
                  \in\operatorname{Br}(W)[n].
\]
The function $t-a$ is a uniformizer at the rational point
$a$, and $\chi$ is unramified there. The cohomological
residue convention of equation (1.17) in the opened
Colliot-Th\'el\`ene--Skorobogatov draft gives
\[
                  \partial_a\beta_{a,\chi}=\chi.
\]
Their Theorem 1.4.10 states explicitly that the Witt residue
is the opposite of this cohomological residue. We use that
sign convention consistently; the purity exact sequence is
unchanged when its residue arrow is negated
\cite[\S1.4.1, equation (1.17), Theorem 1.4.10,
Proposition 3.6.1]{D_ColliotTheleneSkorobogatov2019}.

Every other point of $U$ has valuation zero on $t-a$.
Purity on the regular curve says that a class in
$\operatorname{Br}(W)$ extends to $U$ precisely when its
residue at $a$ is zero. The generic-point injection gives
the first injection of (D.27.2). The explicit symbols give
surjectivity and a splitting on $n$-torsion, proving that
sequence. In particular a nonzero $\beta_{a,\chi}$ cannot
extend across $a$, and adding a constant Brauer class does
not remove its residue.

For ordinary local points avoiding $a$, (D.27.1) is exactly
\[
         \sum_v\operatorname{inv}_v\beta_{a,\chi}(x_v)=0.
\]
A section on $U$ gives global pullbacks and reciprocity for
classes on $U$. The displayed symbols belong to the new
residue summand on $W$, so that assertion does not supply
these tests. A lift or an additional theorem would be
needed. Before even reaching this reciprocity test, one
must prove restricted-product admissibility: D.26 controls
contact with the original boundary, and makes no assertion
that $v(x_v-a)$ is divisible by $n$ at almost every place
for an added interior point $a$.

A small example verifies the distinction between integrality
and the missing reciprocity. This example is not asserted
to arise from a section. Take $k=\mathbf Q$, $a=2$,
$x_v=3$ at every place except $v=3$, and $x_3=4$.
These are ordinary points integral on $W$ outside a finite
set. The translated tuple is $1$ everywhere except at $3$,
where it is $2$. For the quadratic character of
$\mathbf Q(\sqrt3)/\mathbf Q$, the only nonzero local
symbol is $(2,3)_3$: the residue $2$ is a nonsquare modulo
$3$, so its invariant is $1/2$
\cite[example following Corollary 3.6.3]{D_ColliotTheleneSkorobogatov2019}.
All other terms, including the real term, vanish. Thus no
global squareclass has those localizations. The example
need not satisfy the original section's global $t$ and
$1-t$ constraints; it shows precisely that density or even
cofinite integrality alone does not impose the added
reciprocity summand. The exact odd-prime Hilbert-symbol
formula and the residue norm calculation are independently
checked by \path{code/D_canonical_symbol_certificate.py};
its recorded output is
\path{audit/D_canonical_symbol_certificate.json}.

If canonical local matching were proved at every level,
the finite-fibre argument above would give one unique global
completed value. True local norm identities would then give
global norm compatibility, without D.24's failed assertion
that arbitrary choices have lifts. The remaining issue is
the nonemptiness of these canonically constrained finite-level
fibres. A birational section already supplies rational-function
Kummer values, so this criterion does not resolve D2's
stronger cofinite-integrality problem. For a general section
on $U$ it identifies the new residue data needed in the
central part of D4, and does not solve its nonabelian lifting
problem.
\end{proof}
